\subsection{动态加热控制策略优化与数值求解}
在 7.2 建立预冷模型与动态辅助加热控制律后，需要进一步确定参考温升轨迹及反馈增益，使五片电堆在满足启动与安全约束的条件下，以尽可能小的辅助电能完成冷启动。由于功率限幅、状态切换、测量保持和温度过零事件均会改变控制响应，能耗关于控制参数一般不光滑，难以直接采用单一起点的梯度方法求解。因此，本文沿用前两章的“内层瞬态仿真—外层参数优化”框架，在离线仿真中整定控制参数，在线阶段仅根据温度、电压反馈计算分片加热功率，并进一步通过扰动场景筛选获得具有启动裕度的推荐方案。

\subsubsection{控制参数、优化目标与可行域}
动态控制律虽然对五片电池分别计算功率，但各片采用相同的控制结构与参数，功率差异由各片温度、电压、温升速率及估计热负荷决定。因此，无需逐时刻独立优化五条功率曲线，而将常规搜索变量写为
\begin{equation}
\boldsymbol{\theta}=\left(T_{\mathrm{targ}},t_{\mathrm{plan}},K_P,K_I,\Delta T_{\mathrm{taper}},K_r\right).
\tag{7-@@1@@}\label{eq:q4:solver_parameters}
\end{equation}
其中，\(K_r\) 为温升不足补偿增益，其余参数分别控制目标温度、参考升温时间、比例反馈、积分反馈及近目标降功率区间。各参数的搜索范围见表\ref{B7-control-bounds}。\(t_{\mathrm{plan}}\) 是参考轨迹的设计参数，可以超过实际启动期限；候选方案是否可行仍由真实首次成功时刻与实际关热时刻判定，不能以参考轨迹的终点代替启动成功。

\begin{table}[htbp]
    \centering
    \caption{动态辅助加热控制参数的常规搜索范围}
    \label{B7-control-bounds}
    \renewcommand{\arraystretch}{1.5}
    \begin{tabular}{l c c c}
        \toprule
        控制参数 & 单位 & 下界 & 上界 \\
        \midrule
        目标温度 \(T_{\mathrm{targ}}\) & \(^\circ\mathrm C\) & 0.4 & 5 \\
        规划时间 \(t_{\mathrm{plan}}\) & \(\mathrm s\) & 15 & 150 \\
        比例增益 \(K_P\) & \(\mathrm{W\,cm^{-2}\,K^{-1}}\) & 0.015 & 0.5 \\
        积分增益 \(K_I\) & \(\mathrm{W\,cm^{-2}\,K^{-1}\,s^{-1}}\) & 0.001 & 0.08 \\
        渐缩带宽 \(\Delta T_{\mathrm{taper}}\) & \(\mathrm K\) & 0.1 & 2 \\
        温升不足补偿增益 \(K_r\) & \(\mathrm{W\,cm^{-2}}\) & 0 & 0.4 \\
        \bottomrule
    \end{tabular}
\end{table}

为使策略比较具有一致性，采样周期、测量停机规则及安全阈值不参与上述六维搜索。名义搜索中前馈倍率固定为 1，连续保持时间取 \(2\,\mathrm s\)，安全增强起始功率取 \(0.6\,\mathrm{W\,cm^{-2}}\)，冰量预警阈值、电压预警带和综合风险阈值分别取 0.8、\(0.1\,\mathrm V\) 和 0.7，电压创新校正增益取 0.02。另设关闭前馈、比例、积分和温升不足补偿的独立候选，用于检验当前初温下能否仅依靠反应热启动；该候选保留安全后备动作，不受表\ref{B7-control-bounds}中反馈增益正下界的限制。

记满足 7.2 所述功率、真实状态、测量保持与物理有效性条件的参数集合为 \(\Theta_{\mathrm{feas}}\)。首次真实成功需不晚于 \(96.6667\,\mathrm s\)，测量确认关热需在 \(98.6667\,\mathrm s\) 的有限评价窗口内完成。辅助能耗统计至实际关热时刻，归一化目标为
\begin{equation}
\min_{\boldsymbol{\theta}\in\Theta_{\mathrm{feas}}}
J(\boldsymbol{\theta})
=\frac{E_{\mathrm{aux}}(\boldsymbol{\theta})}{E_{\mathrm{ref}}}.
\tag{7-@@2@@}\label{eq:q4:solver_objective}
\end{equation}
式中，\(E_{\mathrm{ref}}\) 取同工况固定恒功率方案的辅助能耗，在同一次搜索中保持不变，仅用于量纲归一化。电压、冰量和内部状态有效性直接作为可行性条件处理；搜索时对不可行候选附加高罚值，但最终选择不允许以能耗收益抵偿安全约束违背。

为避免极小的能耗差掩盖明显的启动时间差，正式精度复算后，在有限可行候选集合 \(\mathcal C_f\) 中先确定最低辅助能耗，再构造其 \(0.05\%\) 近优带：
\begin{equation}
E_{\min}=\min_{\boldsymbol{\theta}\in\mathcal C_f}E_{\mathrm{aux}}(\boldsymbol{\theta}),\qquad
\mathcal C_{\varepsilon}=\left\{\boldsymbol{\theta}\in\mathcal C_f:
E_{\mathrm{aux}}(\boldsymbol{\theta})\leq1.0005E_{\min}\right\}.
\tag{7-@@3@@}\label{eq:q4:solver_energy_band}
\end{equation}
在该集合中，依次优先选择首次真实成功时间较短、最大同刻片间温差较小和辅助能耗较低的候选。用于排序的温差采用式\eqref{eq:q4:comparison_metrics}中的 \(\Delta T_{\max}\)，比较同一时刻五片单电池的温度，能够表征启动过程中的片间一致性；若将不同时刻的最高温度与最低温度相减，则会把整体升温幅度混入空间温差，不能用于本问的一致性评价。

\subsubsection{预冷离散与闭环耦合求解方法}
预冷阶段按照实际层状结构划分 33 个材料区，共设置 146 个控制体，材料界面与控制体边界对齐。设第 \(i\) 个控制体宽度、导热系数及单位面积热容分别为 \(\Delta x_i\)、\(\lambda_i\) 和 \(C_{A,i}\)，则相邻控制体之间及边界控制体至环境之间的单位面积热导为
\begin{equation}
g_{i+1/2}=\left(\frac{\Delta x_i}{2\lambda_i}+\frac{\Delta x_{i+1}}{2\lambda_{i+1}}\right)^{-1},\qquad
g_b=\left(\frac{\Delta x_b}{2\lambda_b}+\frac{1}{h}\right)^{-1}.
\tag{7-@@4@@}\label{eq:q4:solver_precool_conductance}
\end{equation}
边界热导中同时保留半个控制体的导热阻力与外表面对流阻力，使离散散热方向与预冷模型的 Robin 边界一致。令 \(\boldsymbol{u}=\boldsymbol{T}-T_{\mathrm{amb}}\boldsymbol{1}\)，以 \(\mathbf C_{\mathrm p}\) 表示对角热容矩阵、\(\mathbf L_{\mathrm p}\) 表示由内部与边界热导组装的矩阵，则后向欧拉格式为
\begin{equation}
\left(\frac{\mathbf C_{\mathrm p}}{\Delta\tau}+\mathbf L_{\mathrm p}\right)
\boldsymbol{u}^{m+1}=\frac{\mathbf C_{\mathrm p}}{\Delta\tau}\boldsymbol{u}^{m}.
\tag{7-@@5@@}\label{eq:q4:solver_precool_be}
\end{equation}
正式预冷时间步取 \(\Delta\tau=0.25\,\mathrm s\)，并在所需预冷时刻精确落点。得到空间温度场后，按照 7.2 的热容加权映射生成五片单电池和两个端板的初始温度，保留端板与单电池之间的初始热差异。

启动阶段每片采用 58 个基础传输控制体描述水、冰及气体状态，热状态仍由七节点网络计算。为区别于控制采样时刻 \(t_n\)，将内部积分网格记为 \(s_m\)，相邻步长为 \(\Delta s_m=s_{m+1}-s_m\)。给定当前加热功率后，先推进质量迁移与相变，再通过热—电 Picard 迭代求解温度、电压和反应产热，其热网络离散形式为
\begin{equation}
\begin{aligned}
\left(\frac{\mathbf C^{m+1}}{\Delta s_m}+\mathbf G^{m+1}\right)
\boldsymbol{T}^{m+1}
={}&\frac{\mathbf C^{m+1}}{\Delta s_m}\boldsymbol{T}^{m}
+\boldsymbol b_{\mathrm{amb}}^{m+1}+\boldsymbol q_{\mathrm{gen}}^{m+1}\\
&+\frac{\Delta\boldsymbol H_{\mathrm{phase}}^{m+1}}{\Delta s_m}
+\boldsymbol q_{\mathrm{aux}}(s_m).
\end{aligned}
\tag{7-@@6@@}\label{eq:q4:solver_startup_be}
\end{equation}
其中，各热源均按单位面积计量，\(\Delta\boldsymbol H_{\mathrm{phase}}\) 为当前步的相变释热量，辅助源向量的前五个分量为 \(10^4q_k\)，两个端板分量为零。正式积分步长取 \(0.025\,\mathrm s\)，Picard 迭代的温度增量收敛容差取 \(10^{-9}\,\mathrm K\)，每步最多迭代 30 次，并记录实际迭代误差用于最终可行性复核。

控制器每隔 \(0.2\,\mathrm s\) 采集一次温度和电压，经滤波及独立状态观测后更新加热功率；相邻采样点之间功率保持不变。观测器用名义物性、已知电流及已施加的功率独立推进，其内部水冰状态不从被控对象复制。为使数值积分与控制更新一致，积分步在采样时刻、电流平台切换时刻和评价窗口终点截短。步内电流采用累计电荷的精确增量确定，即
\begin{equation}
\bar j_m=\frac{Q(s_{m+1})-Q(s_m)}{\Delta s_m},\qquad
E_{\mathrm{aux}}=25\sum_{s_m<t_{\mathrm{off}}}\Delta s_m\sum_{k=1}^{5}q_k(s_m).
\tag{7-@@7@@}\label{eq:q4:solver_charge_energy}
\end{equation}
式中，\(q_k(s_m)\) 表示当前控制采样区间保持的执行功率，其单位为 \(\mathrm{W\,cm^{-2}}\)，因此面积取 \(25\,\mathrm{cm^2}\) 后能耗单位为焦耳。

当一个积分步跨越五片共同过零事件时，从该步起点重新积分，并以二分法定位首次真实成功时刻；该事件仅用于评价，不提前触发测量停机。停机仍由采样时刻的滤波温度、电压及估计冰量连续满足规定条件 \(2\,\mathrm s\) 后决定，随后永久锁存零功率。每个内部时间步同时检查电压、冰量、剩余气相孔隙率、物种库存、质子电导率及极限电流条件，防止仅凭温度越零或测量保持完成将物理无效轨迹判为成功。

\subsubsection{离线参数搜索与名义节能方案求解}
对于完全冷却、预冷 20 min 和预冷 40 min 三种初态，分别调用相同的闭环仿真器进行参数整定。首先评估关闭名义加热的独立候选。若其实际辅助功率始终为零且满足全部可行性条件，则辅助能耗达到非负下界，无需继续以降低名义辅助能耗为目标搜索。预冷 20 min 工况属于这一情形，但零辅助能耗仅说明该名义初态能够自启动，不能据此推断其具有足够的扰动裕度。

对于需要辅助加热的工况，先对已有物理初值和默认参数进行多起点复核，再以两个固定随机种子运行差分进化搜索。初筛阶段采用 \(0.1\,\mathrm s\) 时间步和每片 29 个传输控制体，以降低大量候选评估的计算成本；随后提取能耗较低且参数互异的 16 个可行候选，统一以 \(0.025\,\mathrm s\) 时间步、每片 58 个控制体重新计算，并以复算后的较优候选为起点进行有界 Powell 局部整定。最终仅在正式精度下的可行候选中，按照式\eqref{eq:q4:solver_energy_band}规定的能耗近优带和排序规则选取名义方案。

上述流程同时覆盖快速升温与延长加热时间两类可能的解，并通过正式精度重排排除粗步长下出现的虚假可行候选。由于实际搜索限于给定反馈结构、参数范围及有限候选集合，所得方案应理解为当前搜索范围内的优选可行解，不能将某次差分进化结果解释为任意连续功率函数空间中的严格全局最优解。

\subsubsection{扰动场景筛选与留裕度控制方案确定}
名义节能方案可能靠近首次成功期限，在初温降低或端部散热增强时失去可行性。因此，在名义参数附近进一步构造留裕度候选：对完全冷却和预冷 40 min 工况，提高目标温度并适当缩短规划时间；对预冷 20 min 工况，恢复小幅比例、积分及温升不足补偿，并比较不同前馈倍率，以获得必要的主动加热储备。

每个候选均在相同的六个训练场景中计算，场景联合考虑初温 \(\pm1\,\mathrm K\) 偏移、片间热导、端板耦合热导及对流系数的 \(\pm20\%\) 偏差，并叠加标准差为 \(0.2\,\mathrm K\) 的温度噪声和 \(0.005\,\mathrm V\) 的电压噪声。除共同物理与停机条件外，训练筛选还要求
\begin{equation}
t_s(\boldsymbol\theta,\omega)\leq94.6667\,\mathrm s,\qquad
t_{\mathrm{off}}(\boldsymbol\theta,\omega)\leq94.6667\,\mathrm s,
\quad \forall\omega\in\Omega_{\mathrm{train}}.
\tag{7-@@8@@}\label{eq:q4:solver_training_deadline}
\end{equation}
该要求相对于首次成功期限保留 \(2\,\mathrm s\) 裕度，相对于允许关热的最终窗口保留 \(4\,\mathrm s\) 裕度。训练初筛采用 \(0.05\,\mathrm s\) 时间步和每片 58 个控制体；从通过全部训练场景的候选中，按名义能耗选取前四个，再以正式步长 \(0.025\,\mathrm s\) 复核名义工况及六个训练场景，最终取复核通过者中名义辅助能耗最低的方案。以下将名义节能方案记为 \(D_{\mathrm N}\)，推荐留裕度方案记为 \(D_{\mathrm R}\)，三种工况的对应参数见表\ref{B7-control-parameters}。

\begin{table}[htbp]
    \centering
    \caption{三种预冷工况的名义节能与推荐留裕度控制参数}
    \label{B7-control-parameters}
    \renewcommand{\arraystretch}{1.5}
    \resizebox{\textwidth}{!}{
    \begin{tabular}{l l c c c c c c c}
        \toprule
        工况 & 方案 & \(T_{\mathrm{targ}}/{}^\circ\mathrm C\) & \(t_{\mathrm{plan}}/\mathrm s\) & \(K_P\) & \(K_I\) & \(\Delta T_{\mathrm{taper}}/\mathrm K\) & \(K_r\) & 前馈倍率 \\
        \midrule
        完全冷却 & \(D_{\mathrm N}\) & 0.45227 & 97.15711 & 0.49466 & 0.02328 & 0.26215 & 0.36643 & 1 \\
        完全冷却 & \(D_{\mathrm R}\) & 3.50000 & 95.21397 & 0.49466 & 0.02328 & 0.26215 & 0.36643 & 1 \\
        预冷 20 min & \(D_{\mathrm N}\) & 1.00000 & 60.00000 & 0 & 0 & 1.00000 & 0 & 0 \\
        预冷 20 min & \(D_{\mathrm R}\) & 1.00000 & 95.00000 & 0.05000 & 0.00500 & 1.00000 & 0.05000 & 0 \\
        预冷 40 min & \(D_{\mathrm N}\) & 0.73357 & 98.57205 & 0.21677 & 0.03518 & 0.60698 & 0.20018 & 1 \\
        预冷 40 min & \(D_{\mathrm R}\) & 2.00000 & 96.60061 & 0.21677 & 0.03518 & 0.60698 & 0.20018 & 1 \\
        \bottomrule
    \end{tabular}}
\end{table}

由表\ref{B7-control-parameters}可见，完全冷却和预冷 40 min 工况主要通过提高目标温度并略微提前参考轨迹增加裕度，而预冷 20 min 工况由零名义加热转为小增益反馈补偿。推荐方案因此承担额外的辅助能耗，其选择目标是在给定训练场景下保留可行性，不应与纯名义条件下的最低能耗混为一谈。

参数确定后予以冻结，再进行独立验证。每种工况在三种初温偏移与十组新噪声种子的组合下开展 30 次试验，三种工况共 90 次；另对三个传热参数分别施加正负偏差，并设置两个联合偏差场景，每种工况 8 次，共 24 次物性失配与噪声联合试验。训练与独立验证使用不同随机种子，后者不再用于逐次调参。相应通过率及失败机理在 7.4 中分析，其结论限于已列出的扰动范围与样本，不能外推为任意扰动下的严格安全保证。

\subsubsection{恒功率对照与预冷时间扫描设置}
为回答动态策略相对于问题三恒功率协同策略的改进效果，首先保留问题三的固定功率分配
\begin{equation}
\overline{\boldsymbol q}_{C}=
(1,\,1,\,0.6245115588,\,1,\,1)\,\mathrm{W\,cm^{-2}}.
\tag{7-@@9@@}\label{eq:q4:solver_constant_power}
\end{equation}
三工况主比较中，将采用测量保持规则的固定恒功率方案记为 \(C_{\mathrm H}\)，其与动态策略使用相同初态、电流曲线、采样滤波和测量保持规则，均以真实首次成功时刻评价启动速度，以实际关热时刻统计辅助电能及启动过程极值。另保留在真实首次成功时立即关热的理想恒功率结果 \(C_0\)，用于复现问题三基准；该结果的统计窗口不同，不与具有测量保持的能耗直接混用。

为区分动态控制收益与重新整定功率分配的收益，进一步对每种初态单独优化恒功率方案，记为 \(C_*\)。先在镜像形式 \((q_e,q_n,q_c,q_n,q_e)\) 下进行差分进化和 Powell 搜索，再加入五片非对称邻域扰动及低功率候选扫描，并以正式精度复算。该对照与动态主方案采用相同保持条件，另外在 \(30\,\mathrm s\)、\(60\,\mathrm s\)、\(90\,\mathrm s\) 及 \(96.6667\,\mathrm s\) 的首次成功期限下构造有限候选能耗—时间前沿。这样可以同时考察固定旧分配、同工况重新优化恒功率与推荐动态控制之间的差异，但恒功率名义再优化结果不代表其经过同等扰动整定后的能力上限。

针对预冷程度的影响，固定采用式\eqref{eq:q4:solver_constant_power}的功率向量，将预冷时间设置为
\begin{equation}
\tau_c\in\{10,15,20,\ldots,100\}\,\mathrm{min}.
\tag{7-@@10@@}\label{eq:q4:solver_cooling_scan}
\end{equation}
对 19 个预冷时刻分别求解空间温度场、映射启动初态并进行恒功率冷启动仿真，记录启动时间、辅助能耗、最低电压、最大冰体积分数及最大同刻片间温差。该扫描沿用问题三的理想首次成功口径，统计窗口截止于 \(t_s\)，不附加测量保持时间；若初态已满足成功条件，则直接记录 \(t_s=0\) 和零辅助能耗。由此可在固定加热方案下单独辨识预冷程度的作用，避免将初始温度变化与控制参数重新整定的影响混合。

最后，对冻结参数后的正式方案分别开展预冷空间与时间加密、启动步长与传输网格加密、控制采样周期变化以及水量和热量守恒检查。三工况主方案还在关热后继续施加题给电流曲线观察 \(60\,\mathrm s\)，该段始终保持零辅助功率，其温度、电压和冰量极值单独统计，用于判断首次启动和短时保持完成后能否继续维持暖态。
